\documentclass[UTF8,12pt]{ctexart}
\usepackage[paperwidth=240mm,paperheight=200mm,margin=14mm,top=8mm]{geometry}
\usepackage{amsmath,amssymb,mathtools}
\usepackage{xcolor}
\pagestyle{empty}
\setlength{\parindent}{0pt}
\setlength{\parskip}{0pt}
\newcommand{\qn}[1]{\textbf{\large #1}\ }
\xeCJKDeclareCharClass{CJK}{"2460 -> "24FF}
\begin{document}
\qn{1992.1}
设 $\begin{cases} x=f(t)-\pi, \\ y=f(\mathrm{e}^{3t}-1), \end{cases}$ 其中 $f$ 可导，且 $f'(0)\neq 0$，则 $\left.\dfrac{\mathrm{d}y}{\mathrm{d}x}\right|_{t=0}=$______.

\vfill
\newpage
\qn{1992.2}
函数 $y=x+2\cos x$ 在 $\left[0,\dfrac{\pi}{2}\right]$ 上的最大值为______.

\vfill
\newpage
\qn{1992.3}
$\lim\limits_{x\to0}\dfrac{1-\sqrt{1-x^2}}{\mathrm{e}^x-\cos x}=$______.

\vfill
\newpage
\qn{1992.4}
$\int_1^{+\infty}\dfrac{\mathrm{d}x}{x(x^2+1)}=$______.

\vfill
\newpage
\qn{1992.5}
由曲线 $y=x\mathrm{e}^x$ 与直线 $y=\mathrm{e}x$ 所围成的图形的面积 $S=$______.

\vfill
\newpage
\qn{1992.6}
当 $x\to0$ 时，$x-\sin x$ 是 $x^2$ 的（　　）
\\
\noindent (A) 低阶无穷小.　　(B) 高阶无穷小.
\\
\noindent (C) 等价无穷小.　　(D) 同阶但非等价的无穷小.

\vfill
\newpage
\qn{1992.7}
设 $f(x)=\begin{cases} x^2, & x\leqslant 0, \\ x^2+x, & x>0, \end{cases}$ 则（　　）
\\
\noindent (A) $f(-x)=\begin{cases} -x^2, & x\leqslant 0, \\ -(x^2+x), & x>0. \end{cases}$
\\
\noindent (B) $f(-x)=\begin{cases} -(x^2+x), & x<0, \\ -x^2, & x\geqslant 0. \end{cases}$
\\
\noindent (C) $f(-x)=\begin{cases} x^2, & x\leqslant 0, \\ x^2-x, & x>0. \end{cases}$
\\
\noindent (D) $f(-x)=\begin{cases} x^2-x, & x<0, \\ x^2, & x\geqslant 0. \end{cases}$

\vfill
\newpage
\qn{1992.8}
当 $x\to1$ 时，函数 $\dfrac{x^2-1}{x-1}\mathrm{e}^{\frac{1}{x-1}}$ 的极限（　　）
\\
\noindent (A) 等于 2.　　(B) 等于 0.
\\
\noindent (C) 为 $\infty$.　　(D) 不存在但不为 $\infty$.

\vfill
\newpage
\qn{1992.9}
设 $f(x)$ 连续，$F(x)=\int_0^{x^2}f(t^2)\,\mathrm{d}t$，则 $F'(x)$ 等于（　　）
\\
\noindent (A) $f(x^4)$.　　(B) $x^2f(x^4)$.
\\
\noindent (C) $2xf(x^4)$.　　(D) $2xf(x^2)$.

\vfill
\newpage
\qn{1992.10}
若 $f(x)$ 的导函数是 $\sin x$，则 $f(x)$ 有一个原函数为（　　）
\\
\noindent (A) $1+\sin x$.　　(B) $1-\sin x$.
\\
\noindent (C) $1+\cos x$.　　(D) $1-\cos x$.

\vfill
\newpage
\qn{1992.11}
求 $\lim\limits_{x\to\infty}\left(\dfrac{3+x}{6+x}\right)^{\frac{x-1}{2}}$.

\vfill
\newpage
\qn{1992.12}
设函数 $y=y(x)$ 由方程 $y-x\mathrm{e}^y=1$ 所确定，求 $\left.\dfrac{\mathrm{d}^2y}{\mathrm{d}x^2}\right|_{x=0}$ 的值.

\vfill
\newpage
\qn{1992.13}
求 $\int\dfrac{x^3}{\sqrt{1+x^2}}\,\mathrm{d}x$.

\vfill
\newpage
\qn{1992.14}
求 $\int_0^\pi\sqrt{1-\sin x}\,\mathrm{d}x$.

\vfill
\newpage
\qn{1992.15}
求微分方程 $(y-x^3)\,\mathrm{d}x-2x\,\mathrm{d}y=0$ 的通解.
设 $f(x)=\begin{cases} 1+x^2, & x\leqslant 0, \\ \mathrm{e}^{-x}, & x>0, \end{cases}$ 求 $\int_1^3 f(x-2)\,\mathrm{d}x$.
求微分方程 $y''-3y'+2y=x\mathrm{e}^x$ 的通解.
计算曲线 $y=\ln(1-x^2)$ 上相应于 $0\leqslant x\leqslant\dfrac{1}{2}$ 的一段弧的长度.
求曲线 $y=\sqrt{x}$ 的一条切线 $l$，使该曲线与切线 $l$ 及直线 $x=0$，$x=2$ 所围成的平面图形面积最小.
已知 $f''(x)<0$，$f(0)=0$，证明对任何 $x_1>0$，$x_2>0$，有 $f(x_1+x_2)<f(x_1)+f(x_2)$.

\vfill
\newpage
\end{document}